\documentclass[book.tex]{subfiles}

\begin{document}

\chapter{Stochastic Gradient Descent}
\label{chap:grad_descent}
\noindent\rule{11cm}{0.4pt} \\
\textbf{Prerequisites:} \autoref{chap:calculus}, \autoref{chap:multidimensional}, \autoref{chap:ml_basics}  \\
\textbf{Difficulty Level:} ** \\
\noindent\rule{11cm}{0.4pt}
\vspace{5mm}

The previous chapters have defined the structure of simple machine learning models but have not explained how to learn model weights from data. The techniques for model learning draw heavily on the optimization and numerical methods we have seen in earlier chapters. The core idea is to use gradient descent to find local optima of loss functions. Consider the simple model 
\begin{align}
h[w, b] = w^T x + b,
\end{align} 
paired with a $L^2$ loss function:

\begin{align}
    \mathcal{L}(h[w, b], x, y) &= | w^T x + b - y| ^2 
\end{align}

Each learning update is to take one step towards a lower loss by following the gradient:

\begin{eqnarray*}
 w &\leftarrow w - \alpha \nabla_w \mathcal{L} \\
 b &\leftarrow b - \alpha \nabla_b \mathcal{L} \\
\end{eqnarray*}

Here the parameter $\alpha$ is called the step size. Let's define this system in Physika

\begin{lstlisting}[language=python]
class LinearModel(w: $\mathbb{R}^m$, b: $\mathbb{R}^m$):
  def $\lambda$(x: $\mathbb{R}^m$) $\to$ $\mathbb{R}$:
    w*x + b
 
def L(model: $[\mathbb{R}^m, \mathbb{R}^m](\mathbb{R}^m\to\mathbb{R})$, x: $\mathbb{R}^m$, y: $\mathbb{R}$) $\to$ $\mathbb{R}$:
  (model(x) - y)$^2$
\end{lstlisting}
The full type signature for \lstinline{LinearModel} can get unwieldy to read at times, so we will often use the convenient shortcut of referring to the type as \lstinline{LinearModel} directly. With this convention, one step of the learning update can then be computed with
\begin{lstlisting}[language=python]
$\alpha$: $\mathbb{R}$
def step(model: LinearModel, x: $\mathbb{R}^m$, y: $\mathbb{R}$):
  model.w = model.w - $\alpha$ * $\nabla$(L(model, x, y), model.w)
  model.b = model.b - $\alpha$ * $\nabla$(L(model, x, y), model.b)
\end{lstlisting}

To learn on a dataset $X \in \mathbb{R}^{N\times m}$, $y \in \mathbb{R}^N$, you can simply repeatedly take gradient descent steps. We  tie this together into a simple learning algorithm for "stochastic gradient descent"

\begin{lstlisting}[language=python]
def sgd(model: $(\mathbb{R}^m\to\mathbb{R})[\mathbb{R}^m, \mathbb{R}^m]$, X: $\mathbb{R}^{N \times m}$, y: $\mathbb{R}^N$, num_epochs: $\mathbb{N}$):
  for _ in num_epochs:
    for i: step(model, $\alpha$, X[i], y[i])
\end{lstlisting}
The code snippet above introduces the new concept of epochs. An epoch is a complete pass of gradient descent steps over the entire dataset. We've tied the implementation above to the particular linear model at hand, but this technique can be applied to arbitrary models. There are a few important variants of this technique. The first is to performed batched updates on a batch of gradients at a time. 
\begin{lstlisting}[language=python]
w: $\mathbb{R}$[m], b: $\mathbb{R}$
def minibatch(X: $\mathbb{R}[B, m]$, y: $\mathbb{R}[B]):$
  $\Sigma_{\nabla w}$, $\Sigma_{\nabla b}$ = zeros(m), zeros(1)
  for i:
    $\Sigma_{\nabla w}$, $\Sigma_{\nabla b}$ += $\nabla$(loss, X[i], y[i], [w, b])
  w = w - $\alpha$ * (1/B)*$\Sigma_{\nabla w}$
  b = b - $\alpha$ * (1/B)*$\Sigma_{\nabla b}$
\end{lstlisting}

In practice, more efficient hardware operations are used than the for-loop above, but conceptually the same operation is performed.

It is an amazing empirical fact that many different architectures and models can be trained with this basic algorithm. One of the biggest mathematical mysteries of deep learning is explaining why gradient descent works in practice. Most loss functions are non-convex, without a theoretically guaranteed unique minimum, but simple gradient descent often does surprisingly well at optimizing these systems. Theorists don't have a definitive understanding of this phenomenon yet.

\section{Exercises}
\begin{enumerate}
    \item Derive and implement the gradient descent step in Physika for a quadratic model
    \begin{align}
        f(x) &= (w^T x + b)^2
    \end{align}
\end{enumerate}
\end{document}